\subsection{Exponential lateral tails in a strip}

Before normalizing the finite chord laws, we control excursions that
travel sideways while remaining in a thin strip. This estimate also
justifies the subsequent lateral exhaustions with a length insertion.

\begin{lemma}[Lateral tail]\label{lem:pv:lateral-tail}
Consider critical exits from a fixed bottom port in a strip of $h$
layers, or in a convex tiling-line truncation of that strip. Let
$X(\gamma)$ be the largest horizontal distance of a visited vertex
from the source. There are absolute constants $C,c>0$ such that
\[
 \sum_{X(\gamma)\ge t}x_*^{|\gamma|}
 \le C\exp\{-ct/(h+1)\}\qquad(t\ge0).
\]
The same bound, with a polynomial prefactor, holds with a first-length
insertion. In particular the contribution with
$X\ge h\log^2h$ to any first-length mass is smaller than every
fixed inverse power of $h$ as $h\to\infty$.
\end{lemma}

\begin{proof}
For the right tail, take two further tiling cuts with outward normals
$(\sqrt3/2,1/2)$ and $(\sqrt3/2,-1/2)$, in coordinates along and
above the bottom side. Let $E_+,E_-$ be the sets of old exits lost when
the corresponding half-plane is imposed. Flux subtraction puts their
phased sums on the distinct rays
$e^{-i\sigma\pi/3}$ and $e^{-2i\sigma\pi/3}$: the turn from the
source to a new cut is fixed. Differently directed cuts do not share
ports. Since vertical coordinates range over $O(h+1)$, their
intercepts can be staggered by $O(h+1)$ and rounded to the tiling grid
so that
\[
 \{X_+>r+C(h+1)\}\subset E_-\subset E_+
 \subset\{X_+>r\},
\]
where $X_+$ is maximum rightward displacement. All old exit weights
have phased real parts comparable to their positive weights by
convexity. Separation of the two rays implies that the absolute
weight of $E_+\setminus E_-$ is at least a fixed positive fraction
of the weight of $E_-$. This gives fixed-factor decay in each
increment $C(h+1)$, starting from the uniformly bounded total exit
mass. Reflection gives the left tail.

In this subtraction, an exit is lost exactly when it visits a removed
vertex. A newly clipped exit attempts such a vertex along the new
cut; an old exit still valid contributes identically to the two flux
sums, even if its port belongs to both boundaries. Thus no boundary
term was omitted by the nesting of the cuts. The constants are
uniform in a convex lateral truncation, and monotone exhaustion
gives the full-strip statement.

Finally a simple path in $h$ layers and horizontal range $X$ visits
at most $C(h+1)(X+1)$ vertices. The tail integral of the preceding
exponential bound consequently gives, after adjusting $C,c$,
\[
 \sum_{X(\gamma)\ge t}|\gamma|x_*^{|\gamma|}
 \le C(h+1)(h+t+1)e^{-ct/(h+1)}.
\]
At $t=h\log^2h$ this is smaller than any fixed power of $h^{-1}$.
\end{proof}

